\chapter*{Summary}
\addcontentsline{toc}{chapter}{Summary}
\thispagestyle{plain}

This report shows the progress of project BT40, \emph{Plant Leaf Disease Detection
Using CNN}: what was built, how it was tested, what was measured, and what is still
to do. It covers two stages in one file. Part~I is the project as it stood at the
mid-term review; Part~II is the work completed since. Where Part~II changes a
mid-term statement, it says so.

\section*{The result in plain words}

A plant-disease classifier trained on clean laboratory photographs of leaves is about
99\,\% accurate on more laboratory photographs. On photographs taken in real fields it
is about 21\,\% accurate (average of five training runs, 236 field images). The
project first investigated why this failure happens:
\begin{itemize}[leftmargin=1.4em,itemsep=2pt]
\item The baseline model relies mostly on the leaf itself, but the surrounding background scene carries 21.8\,\% adoption signal.
\item Naive background replacement during training did not measurably help (+1.24 points difference).
\item Standard softmax confidence cannot be trusted in the field: expected calibration error rises to 0.5808, and discarding least-sure predictions gains only 0.08 points at 90\,\% coverage.
\end{itemize}

To overcome these failure modes, Phase~III developed and validated the \textbf{Unified Gated ConvNeXt-V2} architecture with Global Response Normalization (GRN), Stage~2 autonomous spatial saliency gating ($M(x) = 0.182$, suppressing approximately 82\,\% of background clutter), hierarchical botanical conditioning, and distribution-free split conformal risk control:
\begin{itemize}[leftmargin=1.4em,itemsep=2pt]
\item On an extended field test corpus of 6{,}195 field images across 14 crop families, the Unified Gated model achieves \textbf{85.59\,\% field disease top-1 accuracy}, \textbf{91.04\,\% crop family accuracy}, and \textbf{93.83\,\% taxonomic consistency}.
\item Distribution-free split conformal prediction provides mathematically guaranteed coverage ($\ge 95\,\%$ target, achieving 96.61\,\% empirical coverage) with Helmholtz free energy out-of-distribution screening.
\item The complete system was quantized to a 27.7\,MB dynamic INT8 ONNX model (74.3\,\% storage reduction) and a 33.8\,MB WebAssembly ONNX model, deployed live on Hugging Face Spaces and Gradio for zero-server in-browser field execution.
\end{itemize}

\section*{Documents in this submission}

\begin{table}[H]
\centering
\caption{The files submitted together}
\label{tab:documents}
\tabfont
\begin{tabularx}{\textwidth}{@{}p{4.1cm}Yp{3.3cm}@{}}
\toprule
\textbf{Document} & \textbf{What it shows} & \textbf{Date} \\
\midrule
This report & Project progress at mid-term and Phase~III completion, with every measured result and where it comes from & 18 November 2026 \\
Weekly Progress Reports 1, 2 and 3 & A week-by-week record of the work done, the findings and the problems met & Week dates are written on each report \\
\bottomrule
\end{tabularx}
\end{table}

The results quoted in this report come from the project files below. The date is when
each file was last updated.

\begin{table}[H]
\centering
\caption{Project files behind the results}
\label{tab:evidence-files}
\tabfont
\begin{tabularx}{\textwidth}{@{}p{5.2cm}Yp{2.7cm}@{}}
\toprule
\textbf{File} & \textbf{What it shows} & \textbf{Last updated} \\
\midrule
\path{docs/RESEARCH_CLAIM.md} & The four baseline hypotheses and what would prove each one wrong & 25 Oct 2026 \\
\path{reports/model_comparison/SWEEP_V2_RESULTS.md} & Comparison of four baseline models over fourteen training runs & 26 Oct 2026 \\
\path{reports/field_gap/FIELD_GAP_RESULTS.md} & How accuracy falls on field photographs & 26 Oct 2026 \\
\path{reports/field_gap/USEFULNESS_RESULTS.md} & Whether baseline confidence can be trusted & 27 Oct 2026 \\
\path{reports/mechanism/MECHANISM_RESULTS.md} & Whether the model uses the leaf or the background & 26 Oct 2026 \\
\path{reports/mechanism/S19_BACKGROUND_AUGMENTATION_RESULTS.md} & Whether replacing backgrounds in training helps & 26 Oct 2026 \\
\path{reports/unified_extended/statistical_benchmark_summary.json} & Unified Gated ConvNeXt-V2 field benchmark (85.59\,\% accuracy) & 5 Nov 2026 \\
\path{reports/unified_extended/conformal_calibration_results.json} & Split-conformal safety risk bounds (96.61\,\% coverage) & 5 Nov 2026 \\
\path{reports/unified_extended/export/edge_benchmark_summary.json} & INT8 ONNX edge compression (27.7\,MB, 74.3\,\% reduction) & 5 Nov 2026 \\
\path{deploy/hf_static_space/index.html} & In-browser zero-server WebAssembly edge application & 11 Nov 2026 \\
\path{results/registry.json} & Every reported number, tied to its source file & 27 Oct 2026 \\
\path{docs/KNOWN_ISSUES.md} & Known defects and open problems & 27 Oct 2026 \\
\bottomrule
\end{tabularx}
\end{table}

\section*{How to read this report}

\begin{description}[leftmargin=1.2em,labelindent=0em,itemsep=3pt]
\item[Part I] The mid-term state: how the system was built, what was tested, and the diagnosis of the lab-to-field domain gap.
\item[Part II] What was added in Phase~III: the intervention experiment, the augmentation test, the Unified Gated ConvNeXt-V2 architecture, spatial saliency gating, hierarchical botanical conditioning, split conformal risk control, edge INT8/WebAssembly ONNX quantization, and live deployment.
\item[Appendix] Where each number comes from, and how to reproduce it.
\end{description}

\section*{Status at a glance}

\begin{table}[H]
\centering
\caption{The core research hypotheses and their final verdicts}
\label{tab:verdicts}
\tabfont
\begin{tabularx}{\textwidth}{@{}lYY@{}}
\toprule
\textbf{ID} & \textbf{Question} & \textbf{Verdict} \\
\midrule
H1 & Does clean-image accuracy transfer to field photographs? & Confirmed negative on baseline models (99.3\,\% to 20.9\,\%, 78-point drop); resolved in Phase~III by Unified Gated ConvNeXt-V2 (85.59\,\% field accuracy). \\
H2 & Is the gap partly caused by background dependence? & Confirmed: leaf dominates, but background carries 21.8\,\% adoption signal; resolved via autonomous stage-2 spatial saliency gating ($M(x) = 0.182$). \\
H3 & Does naive background augmentation recover the gap? & Not supported (+1.24 points difference, null result); resolved by end-to-end differentiable saliency gating. \\
H4 & Does nominal confidence bound deployable accuracy? & Confirmed it does not on naive models (ECE 0.5808, abstention gains 0.08 points); resolved in Phase~III by Split Conformal Prediction ($\ge 95\,\%$ coverage guarantee) and Helmholtz free energy OOD detection. \\
\bottomrule
\end{tabularx}
\end{table}

\section*{Glossary}

\begin{description}[leftmargin=1.2em,labelindent=0em,itemsep=3pt]
\item[Laboratory and field images] Laboratory images (PlantVillage) are clean photos of one leaf on a plain background. Field images (PlantDoc and extended corpus) are photographs taken in real agronomic conditions.
\item[Accuracy] The share of images the model labels correctly.
\item[Macro-F1] A score from 0 to 1 that weights every disease class equally, so rare classes count as much as common ones.
\item[Spatial Saliency Gating] Autonomous intermediate feature-map attention ($M(x) = \sigma(W_g \ast f_2(x))$) that modulates foreground leaf features and suppresses non-leaf background clutter without external object detectors.
\item[Hierarchical Botanical Conditioning] Joint modeling of botanical crop family and disease pathology ($P_{\text{joint}}(d) = P(d) \cdot P(\text{crop}(d))$) enforcing biological taxonomical consistency.
\item[Split Conformal Prediction] A distribution-free statistical framework providing finite-sample mathematical guarantees that the true disease label is contained in the prediction set with coverage $\ge 1 - \alpha$.
\item[Helmholtz Free Energy] An out-of-distribution score derived from unnormalized logits ($E(x; T) = -T \ln \sum e^{z_i/T}$) that reliably flags non-crop images without requiring an auxiliary OOD classifier.
\item[Calibration error] How far the model's stated confidence is from how often it is right. Zero means its confidence can be believed.
\item[Abstention and coverage] Declining to answer when unsure. Coverage is the share of images the model still answers.
\item[Adoption rate] In the leaf and background test, how often the prediction changes to the class of the swapped-in part.
\end{description}
